\title{Group Sequence Policy Optimization}

\begin{abstract}
We present Group Sequence Policy Optimization (GSPO), a reinforcement learning algorithm designed for training large language models with high stability, efficiency, and performance. GSPO distinguishes itself from prior approaches by employing sequence-level, rather than token-level, importance ratios; this facilitates clipping, reward computation, and optimization at the sequence granularity. Our results indicate that GSPO not only surpasses the GRPO algorithm in both sample efficiency and training outcomes but also reliably stabilizes reinforcement learning for Mixture-of-Experts (MoE) architectures. Furthermore, GSPO holds promise for streamlining RL infrastructure. These advantages have played a significant role in enhancing the capabilities of the latest Qwen3 models.
\end{abstract}

\section{Introduction}
\label{sec:intro}

Reinforcement learning (RL) has become a central approach for advancing the scalability of language models \citep{o1,dpsk-r1,qwq32b,qwen3}. Utilizing RL at large scales enables language models to engage with complex tasks—including advanced mathematics and programming competitions—by facilitating extended and more intricate reasoning.

For RL scaling to be effective as computational resources increase, it is critical to ensure that training dynamics remain stable and resilient. Despite this necessity, leading RL techniques like GRPO \citep{grpo} suffer from pronounced instability when applied to massive language models, frequently causing irreparable model collapse \citep{qwen3, minimax-m1}. Such instability presents a major obstacle to extending the performance limits of language models through ongoing RL optimization.

This study reveals that GRPO's instability originates from a fundamental flaw in its algorithm: the incorrect application and subsequent invalidation of importance sampling weights. Such misuse leads to substantial variance in training noise, which grows with the length of generated responses and is exacerbated by the algorithm's clipping strategy, ultimately causing model failure.

To overcome these primary issues, we introduce \textbf{Group Sequence Policy Optimization (GSPO)}, a novel reinforcement learning algorithm for large language model training. GSPO's principal contribution is the use of a theoretically justified importance ratio, formulated with respect to sequence likelihood \citep{click}, thereby adhering to the essential tenets of importance sampling. Furthermore, GSPO leverages normalized rewards computed as advantages across multiple responses for each query, promoting congruence between sequence-level reward assignment and the optimization process.

Our experimental results reveal that GSPO markedly outperforms GRPO in terms of training stability, efficiency, and overall effectiveness. Notably, GSPO intrinsically addresses the stability problems commonly encountered during reinforcement learning (RL) of large Mixture-of-Experts (MoE) models, thus obviating the reliance on intricate stabilization mechanisms and pointing towards the possibility of streamlining RL infrastructure. These advantages of GSPO have been instrumental in achieving notable gains in the latest iterations of the Qwen3 models. We anticipate that GSPO will serve as a dependable and extensible foundation for ongoing progress in large-scale RL training of language models.

\section{Preliminaries}

\paragraph{Notation}
Throughout this work, we represent an autoregressive language model with parameters $\theta$ as a policy $\pi_\theta$. Here, $x$ indicates a query, while $\mathcal{D}$ designates the collection of queries. For any response $y$ associated with a given query $x$, the probability of $y$ under the policy $\pi_\theta$ is formulated as $\pi_\theta (y | x)=\prod_{t=1}^{|y|} \pi_\theta (y_t | x, y_{<t} )$, where $|y|$ specifies the total number of tokens in the response $y$. Each query-response tuple $(x, y)$ may be evaluated by a verifier $r$, which produces a reward $r(x, y) \in [0, 1]$.

\paragraph{Proximal Policy Optimization (PPO)} Using samples generated from the old policy $\pi_{\theta_\text{old}}$, PPO \citep{ppo} constrains the policy update within a proximal region of the old policy through the clipping mechanism. Specifically, PPO employs the following objective for policy optimization (we omit the KL regularization term hereinafter for brevity, as it is not the focus of this paper):

\begin{align}
\mathcal{J}_\text{PPO}(\theta) = \mathbb{E}_{ x \sim \mathcal{D},\, y \sim \pi_{\theta_\text{old}}( \cdot | x) }
\left[ \frac{1}{|y|} \sum_{t=1}^{|y|} 
\min \left( w_{t}(\theta) \widehat{A}_{t},  \, \mathrm{clip} \left( w_{t}(\theta), 1 - {\varepsilon}, 1 + {\varepsilon}\right) \widehat{A}_{t} \right)
\right],
\end{align}

Here, the token $y_t$ is assigned an importance ratio given by $
w_{t}(\theta) = \frac{ \pi_{\theta} (y_{t} | x, y_{<t}) }{ \pi_{\theta_\text{old}} (y_{t} | x,y_{<t})} $, and the advantage $\widehat{A}_{t}$ is approximated through a separate value model. The parameter $\varepsilon$ specifies the range used for clipping importance ratios.

A major practical obstacle in PPO arises from its substantial dependence on the value model. Typically, this value model is as large as the policy model, resulting in significant demands on both memory and computational resources. Additionally, the performance of the algorithm is strongly linked to the accuracy of the value estimates it produces. Developing an accurate value model is a difficult task in itself, and making this model suitable for extended sequences and more intricate tasks further amplifies these difficulties.

\paragraph{Group Relative Policy Optimization (GRPO)} GRPO \citep{grpo} bypasses the need for the value model by computing the relative advantage of each response within a group of responses to the same query. Specifically, GRPO optimizes the following objective:

\begin{align}
\label{equ:grpo}
\mathcal{J}_\text{GRPO}(\theta) = \mathbb{E}_{ x \sim \mathcal{D},\, \{y_i\}_{i=1}^G \sim \pi_{\theta_\text{old}}( \cdot | x) }
\left[ \frac{1}{G} \sum_{i=1}^{G} \frac{1}{|y_i|} \sum_{t=1}^{|y_i|} 
\min \left( w_{i,t}(\theta) \widehat{A}_{i,t},  \, \mathrm{clip} \left( w_{i,t}(\theta), 1 - {\varepsilon}, 1 + {\varepsilon}\right) \widehat{A}_{i,t} \right)
\right],
\end{align}

where $G$ is the number of generated responses to each query $x$ (i.e., the group size), and the importance ratio $w_{i,t}(\theta)$ and advantage $\widehat{A}_{i,t}$ of token $y_{i,t}$ are:

\begin{align}
    w_{i,t}(\theta)=\frac{ \pi_{\theta} (y_{i,t} | x, y_{i,<t}) }{ \pi_{\theta_\text{old}} (y_{i,t} | x,y_{i,<t})},\quad\ 
    \widehat{A}_{i,t} = \widehat{A}_{i} = \frac{r(x, y_i) - \mathrm{mean} \left( \{ r(x, y_i) \}_{i=1}^G \right) }{ \mathrm{std} \left( \{ r(x, y_i) \}_{i=1}^G \right) },
\end{align}

respectively, where all the tokens in $y_i$ share the same advantage as $\widehat{A}_{i}$.

\section{Motivation}
\label{sec:motivation}

As models increase in size, incorporate more sparsity (such as in Mixture-of-Experts architectures), and generate longer outputs, maximizing hardware efficiency in RL training requires deploying large rollout batch sizes. To boost sample efficiency, practitioners typically divide this extensive rollout data into several mini-batches for subsequent gradient computations. However, this approach naturally leads to an off-policy learning context, since responses $y$ are drawn according to a previous policy $\pi_{\theta_\text{old}}$ instead of the current policy $\pi_\theta$ targeted during optimization. This scenario underlines why clipping mechanisms are essential in algorithms like PPO and GRPO—they curtail the influence of samples that are substantially “off-policy” during the computation of gradients.

While mechanisms like clipping aim to manage this off-policy discrepancy, we identify a more fundamental issue in GRPO: \textit{its objective is ill-posed}. This problem becomes particularly acute when training large models on long-response tasks, leading to catastrophic model collapse. The ill-posed nature of the GRPO objective stems from a misapplication of importance sampling weights. The principle of importance sampling is to estimate the expectation of a function $f$ under a target distribution $\pi_\text{tar}$ by re-weighting samples drawn from a behavior distribution $\pi_\text{beh}$:

\begin{align}
\mathbb{E}_{z \sim \pi_\text{tar}} \left[ f(z) \right] = \mathbb{E}_{z \sim \pi_\text{beh}} \left[ \frac{ \pi_\text{tar} (z) }{ \pi_\text{beh} (z)} \, f(z) \right].
\end{align}

Fundamentally, this approach depends on taking the average across numerous samples ($N \gg 1$) drawn from the behavior distribution $\pi_\text{beh}$, such that the importance weighting $\frac{ \pi_\text{tar} (z) }{ \pi_\text{beh} (z)}$ can adequately address disparities between the distributions.

By comparison, GRPO utilizes the importance weight $\frac{ \pi_{\theta} (y_{i,t} | x, y_{i,<t}) }{ \pi_{\theta_\text{old}} (y_{i,t} | x,y_{i,<t})}$ at every token position $t$. Given that this weighting is derived from only a single sampled token $y_{i,t}$ for each next-token distribution $\pi_{\theta_\text{old}}( \cdot | x, y_{i, <t})$, it does not accomplish the desired effect of distribution correction. Instead, it injects substantial noise with high variance into the gradients used in training, an issue that compounds over lengthy sequences and becomes worse with the application of clipping. Experimental results show that this can cause the model to collapse—often irreparably. Once collapse sets in, subsequent training fails to recover performance, even after returning to a prior checkpoint, carefully adjusting hyperparameters (such as clipping limits), increasing sequence length, or altering the reinforcement learning queries.

This insight highlights a deep-seated flaw in the GRPO framework. The ineffectiveness of token-level importance weighting underscores a fundamental tenet: \textbf{the level at which the optimization objective operates should align with the level at which rewards are applied}. Since rewards are issued for complete sequences, imposing off-policy corrections at the token level is likely misguided. This leads us to abandon the token-level objective and focus on employing importance weights and optimizing directly at the \textit{sequence level}.

\section{Algorithm}

\subsection{GSPO: Group Sequence Policy Optimization}

Although GRPO faces challenges with the token-level importance weight $\frac{ \pi_{\theta} (y_{i,t} | x, y_{i,<t}) }{ \pi_{\theta_\text{old}} (y_{i,t} | x,y_{i,<t})}$, our findings indicate that, within the framework of language generation, the \textit{sequence-level} importance weight $\frac{ \pi_{\theta} (y | x) }{ \pi_{\theta_\text{old}} (y | x)}$ possesses a well-defined theoretical interpretation: it quantitatively captures the discrepancy between the response $y$ drawn from $\pi_{\theta_\text{old}} (\cdot | x)$ and the current model $\pi_{\theta} (\cdot | x)$. This naturally corresponds to the reward evaluated over entire sequences, and additionally, it can be utilized as an informative guide for implementing the clipping strategy.

Based on this straightforward observation, we propose the \textbf{Group Sequence Policy Optimization (GSPO)} algorithm. GSPO employs the following sequence-level optimization objective:

\begin{align}
\mathcal{J}_\text{GSPO} (\theta) =
\mathbb{E}_{ x \sim \mathcal{D},\, \{y_i\}_{i=1}^G \sim \pi_{\theta_\text{old}}( \cdot | x) }
\left[ 
\frac{1}{G} \sum_{i=1}^{G}
\min \left( s_{i}(\theta)  \widehat{A}_{i},  \, \mathrm{clip} \left( s_{i}(\theta), 1 - {\varepsilon}, 1 + {\varepsilon} \right) \widehat{A}_{i} \right) 
\right],
\label{equ:gspo}
\end{align}

where we adopt the group-based advantage estimation:

\begin{align}
\widehat{A}_{i} = \frac{r(x, y_i) - \mathrm{mean} \left( \{ r(x, y_i) \}_{i=1}^G \right) }{ \mathrm{std} \left( \{ r(x, y_i) \}_{i=1}^G \right) },
\end{align}

and define the importance ratio $s_{i}(\theta)$ based on sequence likelihood \citep{click}:

\begin{align}
s_{i}(\theta) = \left( \frac{ \pi_{\theta} (y_i | x) }{ \pi_{\theta_\text{old}} (y_i | x)} \right)^{\frac{1}{|y_i|}}
=
\exp \left( \frac{1}{|y_i|} \sum_{t=1}^{|y_i|} \log \frac{ \pi_{\theta} (y_{i,t} | x, y_{i,<t}) }{ \pi_{\theta_\text{old}} (y_{i,t} | x,y_{i,<t})} \right).
\end{align}

As a result, GSPO utilizes response-level clipping rather than token-level clipping, in order to remove strongly ``off-policy'' samples from the gradient computation. This approach aligns with both the sequence-based reward structure and the optimization process. Importantly, we implement length normalization for $s_{i}(\theta)$ to mitigate variance and ensure $s_{i}(\theta)$ remains within a consistent numerical interval. Without this normalization, slight changes in token likelihoods may cause large swings in the sequence importance ratio, necessitating different clipping thresholds for responses depending on their lengths. Additionally, it should be noted that the clipping intervals used in GSPO often vary greatly from those in earlier methods (such as GRPO), primarily because the formulations of importance ratios differ substantially.

\subsection{Gradient Analysis}
\label{subsec:gradient}

We can derive the gradient of the GSPO objective as follows (clipping is omitted for brevity):

\begin{align}
\nabla_{\theta} \mathcal{J}_\text{GSPO} (\theta)
=&\ 
\nabla_{\theta} \mathbb{E}_{ x \sim \mathcal{D},\, \{y_i\}_{i=1}^G \sim \pi_{\theta_\text{old}}( \cdot | x) }
\left[ 
\frac{1}{G} \sum_{i=1}^{G} s_{i}(\theta) \widehat{A}_{i}
\right] \\
= &\ 
\mathbb{E}_{ x \sim \mathcal{D},\, \{y_i\}_{i=1}^G \sim \pi_{\theta_\text{old}}( \cdot | x) }
\left[ 
\frac{1}{G} \sum_{i=1}^{G}
s_{i}(\theta) \widehat{A}_{i}
\cdot \nabla_{\theta} \log s_{i}(\theta)
\right] \\
=&\ 
\mathbb{E}_{ x \sim \mathcal{D},\, \{y_i\}_{i=1}^G \sim \pi_{\theta_\text{old}}( \cdot | x) }
\left[ 
\frac{1}{G} \sum_{i=1}^{G} 
\left( \frac{ \pi_{\theta} (y_i | x) }{ \pi_{\theta_\text{old}} (y_i | x)} \right)^{\frac{1}{|y_i|}} \widehat{A}_{i} 
\cdot \frac{1}{|y_i|} \sum_{t=1}^{|y_i|} \nabla_{\theta} \log \pi_{\theta} (y_{i,t} | x, y_{i,<t}) 
\right].
\label{equ:gspo_grad}
\end{align}

For comparison, the gradient of the GRPO objective is as follows (note that $\widehat{A}_{i,t} = \widehat{A}_{i}$):

\begin{align}
\nabla_{\theta} \mathcal{J}_\text{GRPO} (\theta) 
=&\ 
\nabla_{\theta} \mathbb{E}_{ x \sim \mathcal{D},\, \{y_i\}_{i=1}^G \sim \pi_{\theta_\text{old}}( \cdot | x) }
\left[ \frac{1}{G} \sum_{i=1}^{G} \frac{1}{|y_i|} \sum_{t=1}^{|y_i|} 
w_{i,t}(\theta) \widehat{A}_{i,t}
\right] \\
=&\
\mathbb{E}_{ x \sim \mathcal{D},\, \{y_i\}_{i=1}^G \sim \pi_{\theta_\text{old}}( \cdot | x) }
\left[ \frac{1}{G} \sum_{i=1}^{G} \widehat{A}_{i} 
\cdot \frac{1}{|y_i|} \sum_{t=1}^{|y_i|} 
\frac{ \pi_{\theta} (y_{i,t} | x, y_{i,<t}) }{ \pi_{\theta_\text{old}} (y_{i,t} | x,y_{i,<t})} 
\nabla_{\theta} \log \pi_{\theta} (y_{i,t} | x, y_{i,<t})  
\right].
\end{align}

Consequently, the principal difference between GSPO and GRPO stems from \textit{the manner in which gradients of token log likelihoods are weighted}. GRPO assigns a weight to each token based on its corresponding ``importance weight'' $\frac{ \pi_{\theta} (y_{i,t} | x, y_{i,<t}) }{ \pi_{\theta_\text{old}} (y_{i,t} | x,y_{i,<t})}$, resulting in heterogeneous weighting. These variable weights—which can occupy the range $(0, 1+\varepsilon]$ when $\widehat{A}_{i} >0$ or $[1-\varepsilon, +\infty)$ if $\widehat{A}_{i} < 0$—are substantial and their cumulative effects may induce instability or unforeseen behaviors during the learning process. GSPO, on the other hand, treats every token in a generated response uniformly, thereby removing the instability that arises in GRPO from unequal token weighting.

\subsection{GSPO-token: A Token-level Objective Variant}

In scenarios like multi-turn RL, we may desire a finer-grained advantage adjustment than the sequence level. To this end, we introduce a token-level objective variant of GSPO, namely \textbf{GSPO-token}, to allow token-wise advantage customization:

\begin{align}
\mathcal{J}_\text{GSPO-token}(\theta)
=
\mathbb{E}_{ x \sim \mathcal{D},\, \{y_i\}_{i=1}^G \sim \pi_{\theta_\text{old}}( \cdot | x) }
\left[ 
\frac{1}{G} \sum_{i=1}^{G} \frac{1}{|y_i|} \sum_{t=1}^{|y_i|} 
\min \left( s_{i,t}(\theta) \widehat{A}_{i,t},  \, \mathrm{clip} \left( s_{i,t}(\theta), 1 - {\varepsilon}, 1 + {\varepsilon} \right) \widehat{A}_{i,t} \right) 
\right],
\label{equ:gspo-token}
\end{align}

where

\begin{align}
s_{i,t}(\theta) 
= \mathrm{sg} \left[ s_{i}(\theta) \right]  \cdot \frac{ \pi_{\theta} (y_{i,t} | x, y_{i,<t}) }{ \mathrm{sg} \left[ \pi_{\theta} (y_{i,t} | x, y_{i,<t}) \right] },
\end{align}

and $\mathrm{sg}[\cdot]$ denotes only taking the numerical value but stopping the gradient, corresponding to the \texttt{detach} operation in PyTorch. The gradient of GSPO-token can be derived as:

\begin{align}
\nabla_{\theta} \mathcal{J}_\text{GSPO-token}(\theta)
=&\ 
\nabla_{\theta} \mathbb{E}_{ x \sim \mathcal{D},\, \{y_i\}_{i=1}^G \sim \pi_{\theta_\text{old}}( \cdot | x) }
\left[ 
\frac{1}{G} \sum_{i=1}^{G}
\frac{1}{|y_i|} \sum_{t=1}^{|y_i|} 
s_{i,t}(\theta) \widehat{A}_{i,t}
\right] \\
=&\ 
\mathbb{E}_{ x \sim \mathcal{D},\, \{y_i\}_{i=1}^G \sim \pi_{\theta_\text{old}}( \cdot | x) }
\left[ 
\frac{1}{G} \sum_{i=1}^{G} s_i (\theta) \cdot \frac{1}{|y_i|} \sum_{t=1}^{|y_i|} 
\widehat{A}_{i,t} \frac{ \nabla_{\theta} \pi_{\theta} (y_{i,t} | x, y_{i,<t}) }{ \pi_{\theta} (y_{i,t} | x, y_{i,<t}) }
\right] \\
=&\ 
\mathbb{E}_{ x \sim \mathcal{D},\, \{y_i\}_{i=1}^G \sim \pi_{\theta_\text{old}}( \cdot | x) }
\left[ 
\frac{1}{G} \sum_{i=1}^{G}  \left( \frac{ \pi_{\theta} (y_i | x) }{ \pi_{\theta_\text{old}} (y_i | x)} \right)^{\frac{1}{|y_i|}}
\cdot \frac{1}{|y_i|} \sum_{t=1}^{|y_i|}
\widehat{A}_{i,t} \nabla_{\theta} \log \pi_{\theta} (y_{i,t} | x, y_{i,<t})  
\right].
\label{equ:gspo-token_grad}
\end{align}

It should be observed that the expression $\frac{ \pi_{\theta} (y_{i,t} | x, y_{i,<t}) }{ \mathrm{sg} \left[ \pi_{\theta} (y_{i,t} | x, y_{i,<t}) \right] }$ evaluates to 1, thereby ensuring that $s_{i,t}(\theta)$ is equivalent in value to $s_{i}(\theta)$. Examining Equation~\eqref{equ:gspo} alongside Equation~\eqref{equ:gspo-token}, as well as Equation~\eqref{equ:gspo_grad} and Equation~\eqref{equ:gspo-token_grad}, reveals that GSPO-token and GSPO yield matching numerical results for the optimization objective, clipping rule, and theoretical gradient when the advantage is held constant across all tokens in the response $y_i$ (i.e., $\widehat{A}_{i,t} = \widehat{A}_{i}$). Nevertheless, GSPO-token offers greater adaptability, as it allows the advantage to be tailored for each individual token.

\section{Experiments and Discussion}

\subsection{Empirical Results}

We utilize a cold-start model initialized from Qwen3-30B-A3B-Base and evaluate both the training reward progression and model performance trajectories across the AIME'24 benchmark (reporting average Pass@1 from 32 samples), LiveCodeBench (spanning 202410 to 202502, with average Pass@1 over 8 samples), and CodeForces (measured by Elo Rating). Throughout reinforcement learning training, rollout data in each batch is divided into four distinct mini-batches to facilitate gradient updates. For GSPO, the left and right clipping intervals in Equation~\eqref{equ:gspo} are assigned values of 3e-4 and 4e-4, respectively. As a baseline, we employ GRPO, applying left and right clipping boundaries in Equation~\eqref{equ:grpo} at 0.2 and 0.27, based on careful tuning to maintain fairness in comparison. It is important to highlight that GRPO relies on the Routing Replay training method to achieve stable MoE RL convergence, a topic explored further in \S~\ref{subsec:routing_replay}; by contrast, \textit{GSPO removes this requirement}.

As depicted in Figure~\ref{fig:results}, the GSPO training process exhibits stable progress across its duration. Notably, \textbf{GSPO facilitates sustained enhancements in performance by leveraging greater training compute, frequently refreshing the query set, and increasing the output length}. Additionally, GSPO surpasses GRPO in terms of training efficiency, yielding higher training accuracy and superior benchmark outcomes while maintaining parity in training compute and query consumption. Furthermore, GSPO has been effectively utilized for RL training with the latest Qwen3 models, thereby robustly substantiating GSPO's capability to harness RL scaling benefits for large language models.

\subsection{Curious Observation on Clipping Fractions}

GSPO differentiates itself from GRPO by applying clipping to whole sequences, rather than at the level of individual tokens. As illustrated in Figure~\ref{fig:clipping}, this results in a discrepancy of roughly two orders of magnitude between the proportions of clipped tokens in GSPO and GRPO. Notably, even varying the clipping ranges does not significantly affect this substantial gap. Yet, even though GSPO clips far more tokens—meaning fewer tokens are available for model training or gradient calculation—it attains greater training efficiency than GRPO. This surprising outcome, where increasing the fraction of clipped tokens actually boosts training performance, implies that GRPO’s gradient estimates at the token level suffer from intrinsic noise and are suboptimal for exploiting samples. In comparison, the sequence-level strategy employed by GSPO delivers a more dependable and productive learning signal.

\subsection{Benefit of GSPO for MoE Training}
\label{subsec:routing_replay}

\paragraph{Background}
Relative to training dense models with RL methods, the inherently sparse activations found in MoE architectures pose distinctive challenges for training stability. In our investigations, particularly when employing the GRPO algorithm, we observed that the \textit{expert-activation volatility} in MoE models can impede proper convergence during RL optimization. Specifically, after applying one or more steps of gradient-based learning, the selection of experts responsible for generating the same model output may shift considerably. For instance, with the 48-layer Qwen3-30B-A3B-Base model, each RL gradient step results in approximately 10\% of the active experts for a given rollout being different under the updated policy $\pi_{\theta}$ as compared to the previous policy $\pi_{\theta_\text{old}}$.

This instability is exacerbated in MoE models with increased depth, resulting in pronounced oscillations in the token-level importance ratios $w_{i,t}(\theta) = \frac{ \pi_{\theta} (y_{i,t} | x, y_{i,<t}) }{ \pi_{\theta_\text{old}} (y_{i,t} | x,y_{i,<t})}$, which, as elaborated in \S~\ref{sec:motivation} and~\ref{subsec:gradient}, leads to their unreliability and disrupts consistent RL convergence.

\paragraph{Our Previous Approach} In addressing this issue, our earlier solution utilized the \textbf{Routing Replay} training method. This involved storing the experts activated by $\pi_{\theta_\text{old}}$ and subsequently ``replaying'' these routing configurations with $\pi_{\theta}$ when evaluating the importance weights $w_{i,t}(\theta) = \frac{ \pi_{\theta} (y_{i,t} | x, y_{i,<t}) }{ \pi_{\theta_\text{old}} (y_{i,t} | x,y_{i,<t})}$. Through this mechanism, both $\pi_{\theta} (y_{i,t} | x, y_{i,<t})$ and $\pi_{\theta_\text{old}} (y_{i,t} | x,y_{i,<t})$ process $y_{i,t}$ using identical expert activations, thereby maintaining consistency in token-level importance ratios and enabling effective optimization on the same set of activated network paths during gradient updates. As illustrated in Figure~\ref{fig:routing_replay}, Routing Replay is vital for achieving stable convergence in GRPO training within MoE models.

\paragraph{Benefit of GSPO} While Routing Replay assists in stabilizing GRPO training for MoE models, it introduces extra memory and communication burdens by recycling routing configurations, and can impose restrictions on the overall expressive power of the MoE architecture. By contrast, as depicted in Figure~\ref{fig:results}, GSPO completely removes any reliance on Routing Replay, facilitating standard computation of the importance ratios $s_i(\theta)$ and ensuring proper convergence and robust optimization throughout training. The central realization is that GSPO operates solely on the basis of the sequence-level likelihood (i.e., $\pi_{\theta} (y_i | x)$), demonstrating insensitivity to token-wise probabilities (i.e., $\pi_{\theta} (y_{i,t} | x, y_{i,<t})$). Given that the MoE model persistently upholds its language modeling strengths, fluctuations in sequence likelihood remain modest. Overall, GSPO addresses the instability associated with expert selection in MoE models at its root, eliminating the need for workaround methods like Routing Replay. This leads to a more straightforward and stable training regimen and unlocks the model’s full potential, free from imposed limitations.

\subsection{Benefit of GSPO for RL Infrastructure}

Due to inconsistencies in numerical precision between training systems such as Megatron and inference systems like SGLang and vLLM, the standard practice involves recalculating the likelihoods of generated outputs under the previous policy $\pi_{\theta_\text{old}}$ using the training engine. In contrast, GSPO focuses on sequence-level likelihoods during optimization, as opposed to token-level likelihoods; this sequence-level approach is inherently less sensitive to precision mismatches. Consequently, GSPO permits direct utilization of likelihood values produced by the inference engine for optimization purposes, eliminating the necessity for recalculation by the training engine. This capability proves particularly advantageous in contexts such as partial rollouts, multi-turn reinforcement learning, and architectures that separate training and inference components.

\section{Conclusion}

Group Sequence Policy Optimization (GSPO) is introduced as a novel reinforcement learning approach tailored for the training of large language models. Utilizing importance sampling at its core, GSPO constructs importance ratios from sequence likelihoods and implements sequence-level clipping, reward assignment, and optimization procedures. The algorithm surpasses GRPO in terms of training stability, efficiency, and overall performance, especially excelling in large-scale RL scenarios involving MoE models. This groundwork has contributed to the substantial advancements observed in recent Qwen3 models. With GSPO providing a robust and scalable framework, we plan to further expand RL efforts, anticipating major progress in the domain of intelligence.

\end{document}